%% kapa: exam=1 round=22 year=2011 period=2 form=A subject=accounting subject_ko=회계학 seq=22 number=22
%% source: Q-Net 감정평가사 기출문제 (cst00309) artlSeq=5036790 file=extracted/1008/5036790/2011년 제22회 감정평가사 1차시험 기출문제/2교시.hwp
%% conversion: kordoc 4.24.51 -> tools/gen1.py (auto); review=proofread-sonnet; flags=none
\begin{question}{22}
(주)서울은 투자부동산에 대하여는 공정가치모형을, 유형자산에 대하여는 재평가모형을 사용하여 후속측정을 하고 있다. 다음의 자료에 의하여 20×2년에 후속측정과 관련하여 당기손익과 기타포괄손익으로 계상할 금액은?

{\small\centering
\begin{tabular}{|c|c|c|c|}
\hline
구 분 & \makecell{20×1년 초\\취득원가} & \makecell{20×1년 말\\공정가치} & \makecell{20×2년 말\\공정가치} \\
\hline
\makecell{건물\\(투자부동산)} & ￦1,000,000 & ￦1,200,000 & ￦1,100,000 \\
\hline
\makecell{토지\\(유형자산)} & 5,000,000 & 4,750,000 & 5,050,000 \\
\hline
\end{tabular}\par}

\noindent\hspace*{2em}\makebox[3.4cm][l]{\uline{당기손익}}\uline{기타포괄손익}\par
\begin{choices}
  \item \makebox[3.4cm][l]{￦50,000 손실}￦0
  \item \makebox[3.4cm][l]{￦150,000 이익}￦50,000 손실
  \item \makebox[3.4cm][l]{￦150,000 이익}￦50,000 이익
  \item \makebox[3.4cm][l]{￦200,000 이익}￦0
  \item \makebox[3.4cm][l]{￦100,000 이익}￦50,000 손실
\end{choices}
\end{question}
\answer{3}
